% =====================================================================
% Response --- Editor
% ---------------------------------------------------------------------
% The editor's decision letter is not in the repository; the quoted passage
% below is taken from the forwarding message.
% Status: general response written 2026-10-22 (proofread pass: typos,
%         parallelism, nested quotes). Point-level editor items would follow
%         as E2, E3, ... if the letter itself carries further comments.
% =====================================================================
\section*{Editor}

\subsection*{E1 --- Main verdict}
\refcom{[As you will see one report is very negative and one is more
positive. In these circumstance we generally reject the paper. However, in
your case, a member of the Board sent me this message: `I believe it is
possible to provide the authors with a second chance, provided that, in
their revised version, they try to seriously engage and address the major
criticisms raised, especially, by reviewer 2. In conclusion, my summary
recommendation is \emph{Major Revision}, being mindful that this may
ultimately lead to a \emph{Rejection} in view of the very critical comments
received.' On the basis of this evaluation I accepted the idea to give you a
chance.]}
%FIXME [AG, addressed 2026-10-23]: the middle paragraph now leads with the concrete changes (framing as two primitive
% questions -> closure matrix; one open-economy kernel with leaks; two extensions; the appendix apparatus). The apology
% in the first paragraph is left as written -- shortening it is your call, not mine.
\response Many thanks for granting us the opportunity to submit a revised manuscript. We
agree with the reviewers that our original narrative was somewhat ill-focused
-- not only because, given parts of our original framing, it did not spotlight
the most important aspects, but also because it led us to ignore certain
foundational aspects that damaged the credibility of the original text.
Against this backdrop, we decided to rewrite the whole draft, focusing on a
more exhaustive and general treatment of our main object of interest: the
intersection between foundational assumptions of input-output modeling and
more pragmatic applications in the field of transformational studies that rely
on input-output models as a main workhorse. While being more encompassing in
terms of the model space discussed, this fully rewritten new version allows
for discussing our original question in a more useful and better substantiated
way.

Concretely, the revision restates the comparison as two primitive questions --- the
nature of the short run and the question of who bears the cost --- and organises it
as a matrix of five labor-market closures and three financing closures, so that the
fifteen scenarios differ only in the assumptions that define them. The underlying
kernel is stated once, in the open economy, with imported intermediates and product
taxes carried as explicit leaks rather than absorbed into demand, and with the
equivalence results that locate the two poles of the matrix inside one model. Two
extensions probe the limits of the framework: a Kaldor--Verdoorn version in which
capacity responds to demand, and a markup version in which profits enter as a second
income stream. The appendix apparatus documents the kernel, the equivalence proofs,
the calibration and the data pipeline, and the reported tables are generated
directly from the committed run manifests.

In this light, we are extremely thankful for the constructive and direct
feedback received from the reviewers, which guided us in rethinking the main
setup of our narrative.
\loc{The rewrite is the manuscript as a whole; the framing it rests on is stated in Sec.~1 (the two primitive questions and the closure matrix) and Sec.~3 (the family of closures and Table~3), the results in Secs.~5--6.}
% AG (Oct 2026): provisional location --- set the final pointers once the
% manuscript is frozen (the closing-sequence rule: no \loc into a moving text).
